\scanpage{7}
Par réduction au cas de la base $\mathbb Z$ ou $\mathbb Z[1/\unsure{n|2}]$, et par application du Main theorem, on trouve que le plus grand ouvert de ${\unsure{\underline g|\underline{\mathfrak g}}}'$ sur lequel $p$ est un isom est formé des $x'\in\underline g'$ tels que $p^{-1}(x)$ soit fini (ou même, contienne \inserted{un} élément isolé) i.e. l'image inverse de la partie \unsure{ouverte|ouvert} $\underline g^{\mathrm{reg}}\subset\underline g$, formée des points $x$ tels que $q^{-1}(x)$ soit fini (ou comme ...), i.e. des $x\in\underline g$ contenus dans un nb \uline{fini} seulement de sous-algèbres de Borel. En vertu de Springer PM n\textsuperscript{o}30, 5.3, \deleted{\unsure{il semble|il résulte}}\inserted{on voit} que cet ouvert \deleted{\unsure{conti|contient}} rencontre les parties unipotentes de tt\textsuperscript{e} fibre, et on \uline{prévoit} que l'on puisse en déduire qu'il rencontre tte fibre de $f\colon\underline g-\unsure{J|I|\mathfrak J}$.

\uline{Questions}

0) $W$ opère-t-il fidèl\textsuperscript{mt} sur $\underline t^{\unsure{rss|rs}}$ ? \uuline{Non} \uline{en car.} 2, \deleted{\unsure{les|cas}} $A_1,B_2$ ... peut-être d'autres ....

a) $J$ est-il régulier i.e. $J'$ plat sur $J$

b) $f$ est-il plat (impliquant a) ; en fait $f$ plat $\Longleftrightarrow$ $J$ régulier et $f$ équidimensionnel

\margin{OK pour les ``bons'' nombres premiers ?}
\margin{\deleted{\uline{NB} \unsure{la façon d'exprimer} \unsure{ceci} \unsure{l'isom} de l'anneau aff\textsuperscript{ine} de $\underline t/W$ sur \unsure{un} \uline{quotient} \unsure{de} \unsure{le sous-anneau} des invariants de l'anneau affine de $\underline g$ sur $S$, \unsure{$J$} \unsure{est quotient de} \unsure{$f(\underline g)$} ?}}

c) Les fibres de $f$ géom. \unsure{irréductibles|intègres} ? Sont elles chacune l'adhérence d'une orbite sous $G$, le complémentaire d'icelle dans la fibre étant de codim $\geq 2$ (ceci impliquant que les fibres de $f$ sont normales, \unsure{alors} $f$ donc $f'$ normal, donc $\underline g'$ normal \inserted{si base normale}, donc $\underline{\mathscr L}-\underline g'-\underline g$ factorisation de Stein)
\margin{\unsure{*)}}

d) Les points de lissité de $f$ sont ils les points réguliers (quasi-ponctuellement) i.e. ceux au dessus desquels $q$ est \deleted{\unsure{étale} (ou au moins quasi} fini ? Sont-ce aussi les points où le centralisateur \inserted{dans $G$} est de dimension minimum ?, soit $r$ ?

e) \deleted{Les orbites \inserted{unipotentes} de $G$ sur $\tilde G$, et $\underline g$ \unsure{sur}} Définir structure orbitale pour la Chevalley sur $\mathbb Z$, suivant une filtration naturelle de $\underline t$ ; cette structure orbitale devrait être homologue à la structure orbitale de $G$ opérant sur $\tilde G$ (le groupe simplement connexe) ; relations entre les centralisateurs correspondants à des orbites homologues ? / Question analogue pour $\mathscr B,\mathscr L$.

\margin{f) Reprendre théorie de Kostant : A-t-on section de l'ouvert des pts réguliers etc ---}

\margin{g) Étudier fibre de $p\colon\underline{\mathscr L}-\underline g'$, i.e. les composantes connexes de la variété des sous-alg. de Borel de $\underline g$ contenant un élément donné. Quelle est sa dimension (relation avec dim du \unsure{stabilisateur|fixeur|fixateur|l'orbite} )?}
% archive sheet number: 6
